\newpage

\section*{Anotácia}

Slovenská technická univerzita v Bratislave \\
\myFacultySK \\[4pt]
Študijný program: \myStudyProgramSK \\
Autor: \myName \\
Diplomová práca: \myTitleSK \\
Vedúci záverečnej práce: \mySupervisor \\
\myDateSK

\vspace{8pt}

Empirické zistenia o vplyve nástrojov umelej inteligencie na učenie si protirečia: adaptívny tutor prekonáva aktívne učenie, generatívny nástroj však znižuje výsledky na následnom teste a priame porovnanie adaptívneho a generického nástroja nenachádza rozdiel. Práca tvrdí, že rozpor je zdanlivý, a zjednocuje ho hypotézou, podľa ktorej nástroj zvyšuje výkon počas hodiny a pocit naučenia takmer vždy, kým jeho vplyv na trvalé osvojenie závisí od toho, či interakcia žiaka s ním zostáva generatívna, alebo sa stáva offloadingová. Táto veličina dosiaľ nebola meraná priamo. Práca ju operacionalizuje ako index epistemickej kvality promptov~--- taxonómiu žiackych dopytov ukotvenú v rámci ICAP spolu s automatickým klasifikátorom~--- a počíta ju z promptov zachytených nástrojom vyvinutým v technickej časti práce. Druhým príspevkom je návrh krížového experimentu na slovenských stredných školách v usporiadaní replikovaného latinského štvorca so šiestimi sekvenciami, ktorý oddeľuje účinok pedagogického ohraničenia nástroja od účinku modelu tým, že generickú a adaptívnu podmienku stavia nad ten istý model, tie isté úlohy a to isté rozhranie. Merané je okamžité aj odložené vybavenie, transfer, kalibrácia sebadôvery a procesné ukazovatele, pričom obe testové formy sú prepojené kotviacimi položkami a kalibrované súbežne. Práca dôsledne rozlišuje, čo randomizácia identifikuje a čo nie: do konfirmačnej rodiny vstupujú len tvrdenia identifikované dizajnom, kým mechanizmus sa testuje trianguláciou troch odhadov pod tromi rôznymi predpokladmi a analýzou citlivosti na nepozorované zmätenie. Analytický plán je uzavretý, implementovaný a pred zberom overený na dátach zo známeho generujúceho modelu; validácia odhalila šesť chýb návrhu a viedla k zvýšeniu cieľovej veľkosti vzorky z 12 na 18 tried.

\vspace{6pt}
\textbf{Oblasť problematiky:} umelá inteligencia vo vzdelávaní, kognitívny offloading, krížový experiment, metakognitívna kalibrácia, analýza promptov


%%%%%%%%%%%%%%%%%%%%%%%%%%%%%%%%%%%%%%%%%%%%%
\vfill
%%%%%%%%%%%%%%%%%%%%%%%%%%%%%%%%%%%%%%%%%%%%%
\section*{Annotation}

Slovak University of Technology in Bratislava \\
\myFaculty \\[4pt]
Study programme: \myStudyProgram \\
Author: \myName \\
Master's Thesis: \myTitle \\
Supervisor: \mySupervisor \\
\myDate

\vspace{8pt}

Empirical findings on how artificial intelligence tools affect learning appear to contradict one another: an adaptive tutor outperforms in-class active learning, a generative tool lowers subsequent test scores, and a direct comparison of an adaptive with a generic tool finds no difference. This thesis argues that the contradiction is only apparent and resolves it with a single hypothesis: a tool raises in-class performance and the feeling of learning almost unconditionally, whereas its effect on \emph{durable} learning depends on whether the student's interaction with it remains generative or becomes offloading. That quantity has never been measured directly. The thesis operationalises it as the Prompt Epistemic Quality Index~--- a taxonomy of student prompts anchored in the ICAP framework together with an automatic classifier~--- computed from prompts captured by the tool developed in the engineering part of the work. The second contribution is the design of a cross-over classroom experiment in Slovak secondary schools, arranged as a replicated Latin square with six sequences, which separates the effect of the tool's pedagogical scaffolding from the effect of the underlying model by building the generic and the adaptive condition on the same model, the same exercises and the same interface. Immediate and delayed retention, far transfer, confidence calibration and process measures are recorded, with the two test forms linked by common anchor items and calibrated concurrently. The thesis separates what randomisation identifies from what it does not: only design-identified claims enter the confirmatory family, while the mechanism is tested by triangulating three estimators resting on three different assumption sets, together with a sensitivity analysis for unmeasured confounding. The analysis plan is frozen, implemented and validated before data collection on data from a known generating model; the validation exposed six design faults and raised the target sample from 12 to 18 classes.

\vspace{6pt}
\textbf{Keywords:} artificial intelligence in education, cognitive offloading, cross-over design, metacognitive calibration, prompt analysis
